\chapter{Introduction}

\section{研究背景}

在碳中和政策与城市绿色物流发展的推动下，新能源物流车已广泛应用于城市配送环节。与传统燃油车不同，电动物流车受电池容量与载重上限的双重约束，同时还面临续航焦虑、途中补能需求以及充电站排队与负荷压力等问题。城市配送任务往往随时间动态到达，任务地点与货物重量具有随机性，若仅依赖人工排班，很难在时效性、行驶成本与能耗之间取得稳定平衡。因此，从中央仓库管理者的视角出发，对有限规模的新能源车队进行协同调度与路径规划，已成为城市末端物流运营中的关键课题。

\section{问题与挑战}

本课题所关注的核心场景可概括为：仓库配置一支车型各异、电量与载重能力不同的新能源车队，在基于图结构的道路网络上承接动态出现的配送任务。每辆车在行驶过程中需同时满足路网连通性、剩余电量、当前载重以及充电站容量与排队等约束；当电量不足以抵达下一服务点或返回仓库时，还需在多个充电站之间做出补能决策。任务侧则包含产生时间、空间坐标、货物重量与完成时限等属性，完成越早、路径越短通常收益越高，超时则产生惩罚。

在上述约束下，调度与路径规划问题本身具有高度的组合复杂性。不同的调度思想——例如最近任务优先、最大载重优先、截止时间优先、综合评分策略，乃至静态全局优化与元启发式方法——在数据结构设计、状态更新、充电预判与评分口径上往往需要各自独立的实现逻辑。若缺少统一的仿真环境与可比的评价指标，算法之间的优劣难以在同一问题规模下客观衡量，开发与验证成本也会显著上升。

\section{项目动机}

新能源物流车队协同调度已在产业实践中逐步落地，但面向具体场景的调度算法极其多样，实现细节差异大、调试周期长。为应对这一问题，本项目并不将某一种算法固化在孤立脚本中，而是构建一套可扩展的仿真框架：在统一的任务生成、车辆演化、充电站排队与评分机制下，接入多种动态调度策略，并可选地与静态全局优化结果进行对比。通过仿真复现“中央仓库管理者”的决策过程，能够在小、中、大等不同规模配置下观察各策略在完成率、总分、路径长度与能耗等指标上的差异，为算法选型与后续改进提供可重复、可量化的实验依据。

\section{本项目工作}

本项目实现了 NEFT（New Energy Fleet Transportation）新能源物流车队协同调度仿真系统。系统以图结构描述城市道路网络，支持基于真实地图数据的路径计算；在仿真过程中动态生成任务，跟踪每辆车的位置、电量、载重与任务状态，并在电量不足时自动选择兼顾距离、负荷与排队的充电站。调度层提供多种实时策略（如最近任务优先、最大载重优先、优先级与综合评分策略等），并统一嵌入电量预判与充电改派逻辑，保证不同算法在相同物理约束下公平比较。此外，系统支持静态规划模式，可利用 OR-Tools 或 MIP 类求解器在任务已知的前提下求取参考性较强的全局方案，与动态策略结果对照分析差距及原因。前端通过 Web 界面与地图可视化展示车队运行过程，便于演示与结果汇报。

\section{报告结构}

本报告其余章节安排如下：第二章介绍系统总体架构与主要模块划分；第三章阐述道路网络、任务、车辆与充电站等核心数据模型；第四章说明动态调度策略与静态优化方法的设计与实现；第五章给出不同规模与策略下的仿真实验与对比分析；第六章总结全文并讨论后续改进方向。
